% 表 1.4 Jones 忠实表示符号
\noindent\resizebox{\textwidth}{!}{%
\begin{tabular}{|c|c|c|c|c|c|c|c|}
\hline
\multicolumn{8}{|c|}{(a) \textit{立方}} \\ \hline
$E$ & $xyz$ & $I$ & $\bar{x}\bar{y}\bar{z}$ & $C_{4x}^{+}$ & $x\bar{z}y$ & $S_{4x}^{-}$ & $\bar{x}z\bar{y}$ \\ \hline
$C_{2x}$ & $x\bar{y}\bar{z}$ & $\sigma_{x}$ & $\bar{x}yz$ & $C_{4y}^{+}$ & $zy\bar{x}$ & $S_{4y}^{-}$ & $\bar{z}yx$ \\ \hline
$C_{2y}$ & $\bar{x}y\bar{z}$ & $\sigma_{y}$ & $x\bar{y}z$ & $C_{4z}^{+}$ & $\bar{y}xz$ & $S_{4z}^{-}$ & $y\bar{x}\bar{z}$ \\ \hline
$C_{2z}$ & $\bar{x}\bar{y}z$ & $\sigma_{z}$ & $xy\bar{z}$ & $C_{4x}^{-}$ & $xz\bar{y}$ & $S_{4x}^{+}$ & $\bar{x}\bar{z}y$ \\ \hline
$C_{31}^{+}$ & $zxy$ & $S_{61}^{-}$ & $\bar{z}\bar{x}\bar{y}$ & $C_{4y}^{-}$ & $\bar{z}yx$ & $S_{4y}^{+}$ & $z\bar{y}\bar{x}$ \\ \hline
$C_{32}^{+}$ & $\bar{z}x\bar{y}$ & $S_{62}^{-}$ & $z\bar{x}y$ & $C_{4z}^{-}$ & $y\bar{x}z$ & $S_{4z}^{+}$ & $\bar{y}x\bar{z}$ \\ \hline
$C_{33}^{+}$ & $\bar{z}\bar{x}y$ & $S_{63}^{-}$ & $zx\bar{y}$ & $C_{2a}$ & $yx\bar{z}$ & $\sigma_{da}$ & $\bar{y}\bar{x}z$ \\ \hline
$C_{34}^{+}$ & $z\bar{x}\bar{y}$ & $S_{64}^{-}$ & $\bar{z}xy$ & $C_{2b}$ & $\bar{y}\bar{x}\bar{z}$ & $\sigma_{db}$ & $yxz$ \\ \hline
$C_{31}^{-}$ & $yzx$ & $S_{61}^{+}$ & $\bar{y}\bar{z}\bar{x}$ & $C_{2c}$ & $z\bar{y}x$ & $\sigma_{dc}$ & $\bar{z}y\bar{x}$ \\ \hline
$C_{32}^{-}$ & $y\bar{z}\bar{x}$ & $S_{62}^{+}$ & $\bar{y}zx$ & $C_{2d}$ & $\bar{z}\bar{y}\bar{x}$ & $\sigma_{dd}$ & $zyx$ \\ \hline
$C_{33}^{-}$ & $\bar{y}z\bar{x}$ & $S_{63}^{+}$ & $y\bar{z}x$ & $C_{2e}$ & $\bar{x}zy$ & $\sigma_{de}$ & $x\bar{z}\bar{y}$ \\ \hline
$C_{34}^{-}$ & $\bar{y}\bar{z}x$ & $S_{64}^{+}$ & $yzx$ & $C_{2f}$ & $\bar{x}\bar{z}\bar{y}$ & $\sigma_{df}$ & $xzy$ \\ \hline
\end{tabular}}

\smallskip

\noindent\resizebox{0.62\textwidth}{!}{%
\begin{tabular}{|c|l|c|l|}
\hline
\multicolumn{4}{|c|}{(b) \textit{六方}} \\ \hline
$E$ & $x, y, z$ & $I$ & $\bar{x}, \bar{y}, \bar{z}$ \\ \hline
$C_{6}^{+}$ & $x-y, x, z$ & $S_{3}^{-}$ & $y-x, \bar{x}, \bar{z}$ \\ \hline
$C_{3}^{+}$ & $\bar{y}, x-y, z$ & $S_{6}^{-}$ & $y, y-x, \bar{z}$ \\ \hline
$C_{2}$ & $\bar{x}, \bar{y}, z$ & $\sigma_{h}$ & $x, y, \bar{z}$ \\ \hline
$C_{3}^{-}$ & $y-x, \bar{x}, z$ & $S_{6}^{+}$ & $x-y, x, \bar{z}$ \\ \hline
$C_{6}^{-}$ & $y, y-x, z$ & $S_{3}^{+}$ & $\bar{y}, x-y, \bar{z}$ \\ \hline
$C'_{21}$ & $y-x, y, \bar{z}$ & $\sigma_{d1}$ & $x-y, \bar{y}, z$ \\ \hline
$C'_{22}$ & $x, x-y, \bar{z}$ & $\sigma_{d2}$ & $\bar{x}, y-x, z$ \\ \hline
$C'_{23}$ & $\bar{y}, \bar{x}, \bar{z}$ & $\sigma_{d3}$ & $y, x, z$ \\ \hline
$C''_{21}$ & $x-y, \bar{y}, \bar{z}$ & $\sigma_{v1}$ & $y-x, y, z$ \\ \hline
$C''_{22}$ & $\bar{x}, y-x, \bar{z}$ & $\sigma_{v2}$ & $x, x-y, z$ \\ \hline
$C''_{23}$ & $y, x, \bar{z}$ & $\sigma_{v3}$ & $\bar{y}, \bar{x}, z$ \\ \hline
\end{tabular}}

\smallskip
\parbox{0.92\textwidth}{\textit{表 1.4 的注.}
(i) 如正文所述，(主动) 算符 $R$ 旁边的符号就是 $R$ 作用在向量 $(x, y, z)$ 上得到的向量。
(ii) 在 (a) 中单位矢量 $\mathbf{i}$、$\mathbf{j}$、$\mathbf{k}$ 互相垂直，沿图 1.1 的 $x$、$y$、$z$ 轴；在 (b) 中单位矢量 $\mathbf{i}$、$\mathbf{j}$、$\mathbf{k}$ 沿图 1.2 的 $x$、$y$、$z$ 轴，因此 $\mathbf{i}$ 与 $\mathbf{j}$ 互成 120° 且都在垂直于 $\mathbf{k}$ 的平面内。
(iii) 例如 $C_{4z}^{-}$ 的符号为 $(y\bar{x}z)$，于是 $C_{4z}^{-}(x\mathbf{i} + y\mathbf{j} + z\mathbf{k}) = (y\mathbf{i} - x\mathbf{j} + z\mathbf{k})$。\footnote{扫描件此处印作 $(\bar{y}\bar{x}z)$，与其下方自带公式 $(y\mathbf{i} - x\mathbf{j} + z\mathbf{k})$ 即 $(y\bar{x}z)$ 矛盾，按公式取 $(y\bar{x}z)$。——中译者}
(iv) (a) 亦可用于四方、正交、单斜、三斜晶系；对单斜与三斜晶系，$\mathbf{i}$、$\mathbf{j}$、$\mathbf{k}$ 平行于晶体坐标轴。(b) 亦可用于三方晶系。
(v) (a) 与 (b) 中的符号其实就是《X 射线晶体学国际表》（Henry and Lonsdale 1965）第一卷表 4.3 中点式立方与六方空间群的一般等价位置坐标。}
